\section{Incentivized Curation Market}

\label{sec:curation}

\subsection{Market Design}

The curation market rewards participants for improving data quality through a bonding curve mechanism:

\begin{equation}
    \text{Reward}(\text{action}, \text{ZDO}) = R_{\text{base}}(\text{action}) \cdot \Delta Q(\text{ZDO}) \cdot B(\text{ZDO})
\end{equation}

\noindent where $R_{\text{base}}$ is the base reward for the action type, $\Delta Q$ is the quality improvement, and $B$ is a bonding curve that increases rewards for curating underserved taxonomic groups and geographic regions.

\begin{table}[H]
\centering
\caption{Curation actions and base reward ranges.}
\label{tab:curation_rewards}
\begin{tabularx}{\textwidth}{@{}lcX@{}}
\toprule
\textbf{Action} & \textbf{Base Reward (ZOO)} & \textbf{Description} \\
\midrule
Contribution & 10--500 & New dataset ingestion with metadata \\
Validation & 5--50 & Independent quality verification \\
Annotation & 2--20 & Taxonomic or geographic annotation refinement \\
Translation & 5--30 & Metadata translation to additional languages \\
Integration & 20--200 & Cross-dataset linkage and harmonization \\
\bottomrule
\end{tabularx}
\end{table}

\subsection{Underserved Region Bonding Curve}

To combat geographic and taxonomic bias, the bonding curve $B$ provides multiplied rewards for data from underrepresented areas:

\begin{equation}
    B(\text{ZDO}) = 1 + \beta_{\text{geo}} \cdot \left(1 - \frac{|\mathcal{D}_{\text{cell}}|}{|\mathcal{D}_{\text{median}}|}\right)^+ + \beta_{\text{taxon}} \cdot \left(1 - \frac{|\mathcal{D}_{\text{taxon}}|}{|\mathcal{D}_{\text{median\_taxon}}|}\right)^+
\end{equation}

\noindent where $|\mathcal{D}_{\text{cell}}|$ is the number of existing records in the geographic grid cell containing the ZDO, $|\mathcal{D}_{\text{median}}|$ is the median records per cell globally, and analogous terms apply for taxonomic coverage.
The $(\cdot)^+$ operator ensures the bonus is non-negative.

\subsection{Validation Protocol}

Data validation follows a stake-weighted consensus mechanism:

\begin{enumerate}[leftmargin=*]
    \item A contributor submits a ZDO with metadata and stakes $s_c$ tokens.
    \item The ZDO is assigned to $k$ randomly selected validators (minimum 3).
    \item Each validator independently assesses quality scores and stakes $s_v$ tokens on their assessment.
    \item Quality scores are aggregated using weighted median (weights proportional to validator reputation).
    \item Validators whose scores are within one standard deviation of the consensus receive rewards; outliers lose stakes.
\end{enumerate}

This mechanism is incentive-compatible: honest assessment maximizes expected payoff for validators with accurate domain knowledge.

% ============================================================
% 7. Implementation
% ============================================================
